% Modul 8 (UTS), dihasilkan oleh tools/build.py
\documentclass[a4paper]{article}
\usepackage{modulITERA}
\begin{document}
\Sampul{8}{Ujian Tengah Semester}{Kisi-kisi, format, tata tertib, dan contoh soal untuk materi Pertemuan 1 sampai 6}{Bagian A untuk CPMK0607 dan Bagian B untuk CPMK0608, masing-masing 50 poin}{120 menit ujian, ditambah latihan mandiri secukupnya}{\texttt{02 Hands-on/P8 - UTS - Contoh Soal - Hands-on/}}
\Bagian{Pembuka}{Tentang modul ini}
Modul ini membantu kalian bersiap menghadapi UTS Praktikum Pemrograman. Isinya format ujian, kisi-kisi per topik, tata tertib, empat contoh soal yang setara dengan soal ujian, dan rubrik penilaian. Contoh soal ada di folder hands-on P8 dan bisa dicek dengan \texttt{cek.sh}.

\Bagian{Ujian}{Format UTS}
\par\vspace{4pt}\noindent{\fontsize{9.5pt}{13pt}\selectfont
\begin{tabular}{@{}>{\raggedright\arraybackslash}p{0.080\textwidth}>{\raggedright\arraybackslash}p{0.500\textwidth}>{\raggedright\arraybackslash}p{0.120\textwidth}>{\raggedright\arraybackslash}p{0.080\textwidth}>{\raggedright\arraybackslash}p{0.200\textwidth}@{}}
\kepala{Soal} & \kepala{Bentuk} & \kepala{Bagian} & \kepala{Poin} & \kepala{CPMK}\\
1 & Menulis program: masukan, ekspresi, dan percabangan & A & 25 & CPMK0607\\\garis
2 & Menulis program: perulangan dengan sentinel atau pola & A & 25 & CPMK0607\\\garis
3 & Tracing program dengan percabangan dan perulangan & B & 25 & CPMK0608\\\garis
4 & Menemukan, menjelaskan, dan memperbaiki kesalahan & B & 25 & CPMK0608\\
\garisdasar
\end{tabular}}\par\vspace{6pt}
\Bagian{Ujian}{Kisi-kisi}
\par\vspace{4pt}\noindent{\fontsize{9.5pt}{13pt}\selectfont
\begin{tabular}{@{}>{\raggedright\arraybackslash}p{0.240\textwidth}>{\raggedright\arraybackslash}p{0.130\textwidth}>{\raggedright\arraybackslash}p{0.490\textwidth}>{\raggedright\arraybackslash}p{0.120\textwidth}@{}}
\kepala{Topik} & \kepala{Pertemuan} & \kepala{Indikator yang diuji} & \kepala{Bagian}\\
Struktur program dan keluaran & P1 & menulis keluaran terformat dengan escape sequence & A\\\garis
Variabel, tipe, masukan & P2 & memilih tipe, membaca masukan, pembagian bulat & A, B\\\garis
Operator dan precedence & P3 & \texttt{/} dan \texttt{\%}, urutan evaluasi, ekspresi logika & A, B\\\garis
Percabangan & P4 & rantai \texttt{else if}, \texttt{switch}, nilai batas & A, B\\\garis
Perulangan & P5 & \texttt{for}, \texttt{while}, sentinel, akumulator & A, B\\\garis
Perulangan bersarang & P6 & pola, \texttt{break}, \texttt{continue}, tracing bersarang & A, B\\
\garisdasar
\end{tabular}}\par\vspace{6pt}
\Bagian{Ujian}{Tata tertib}
\begin{Jawaban}[Yang boleh]
Satu lembar catatan A4 tulisan tangan sendiri (dua sisi). GDB Online atau g++ di komputer lab untuk soal A.
\end{Jawaban}
\begin{Tebak}[Yang tidak boleh]
AI, internet selain GDB Online, catatan elektronik, berdiskusi, atau membuka file selain lembar jawaban.
\end{Tebak}
\begin{Penting}[]
Soal B dikerjakan di kertas tanpa menjalankan program. Pelanggaran tata tertib membuat nilai UTS 0 dan dilaporkan ke program studi.
\end{Penting}
\Bagian{Latihan}{Contoh soal}
Contoh soal ada di \texttt{02 Hands-on/P8 - UTS - Contoh Soal - Hands-on/latihan/}. Kerjakan dengan batas waktu seperti ujian, lalu cek dengan \texttt{bash cek.sh}. Kunci jawaban disimpan dosen.
\Sub{Soal 1 (Bagian A): masukan, pembulatan, \texttt{switch}, potongan}
Hitung ongkos kirim dari berat (gram) dan zona dengan pembulatan ke atas per kilogram dan potongan 10\% untuk paket di atas 10 kg.

\begin{MKeluaran}[title={Contoh masukan}]
1001 A
12500 C
300 D
\end{MKeluaran}
\begin{MKeluaran}[title={Keluaran yang diharapkan}]
2 kg, ongkos 18000
13 kg, ongkos 175500
Zona tidak valid
\end{MKeluaran}
\Sub{Soal 2 (Bagian A): sentinel, akumulator, maks-min, dua desimal}
Baca suhu harian sampai 999, lalu tampilkan banyak hari, rata-rata, banyak hari panas, dan selisih tertinggi-terendah.

\begin{MKeluaran}[title={Contoh masukan}]
30 33 29 35 31 999
\end{MKeluaran}
\begin{MKeluaran}[title={Keluaran yang diharapkan}]
Hari      : 5
Rata-rata : 31.60
Hari panas: 2
Selisih   : 6
\end{MKeluaran}
\Sub{Soal 3 (Bagian B): trace table dengan \texttt{if-else} dan \texttt{break}}
Buat trace table untuk \texttt{x}, \texttt{y}, dan \texttt{i}, lalu tuliskan keluaran program tanpa menjalankannya.

\begin{MKeluaran}[title={Keluaran yang diharapkan}]
1: 2 18
2: 4 18
3: 4 14
akhir 16 14
\end{MKeluaran}
\Sub{Soal 4 (Bagian B): temukan tiga kesalahan, jelaskan akibat, perbaiki}
Program rata-rata nilai lulus memberi hasil salah. Sebutkan tiga kesalahan dengan nomor barisnya, jelaskan akibatnya, dan tulis versi yang benar.

\begin{MKeluaran}[title={Contoh masukan}]
80 55 70 60 45
\end{MKeluaran}
\begin{MKeluaran}[title={Keluaran yang diharapkan}]
Rata-rata lulus: 70
\end{MKeluaran}
\Bagian{Penilaian}{Rubrik penilaian ujian}
\par\vspace{4pt}\noindent{\fontsize{8.5pt}{11.5pt}\selectfont
\begin{tabular}{@{}>{\raggedright\arraybackslash}p{0.180\textwidth}>{\raggedright\arraybackslash}p{0.240\textwidth}>{\raggedright\arraybackslash}p{0.200\textwidth}>{\raggedright\arraybackslash}p{0.200\textwidth}>{\raggedright\arraybackslash}p{0.160\textwidth}@{}}
\kepala{Aspek} & \kepala{Sangat baik 85–100} & \kepala{Baik 70–84} & \kepala{Cukup 55–69} & \kepala{Kurang <55}\\
A. Program berjalan & kompilasi bersih, keluaran tepat untuk semua uji & keluaran tepat untuk sebagian besar uji & berjalan, keluaran sering salah & tidak terkompilasi\\\garis
A. Struktur kode & tipe, nama variabel, dan alur rapi & satu kekurangan kecil & sulit dibaca & tidak sesuai soal\\\garis
B. Tracing & semua nilai dan keluaran tepat & satu langkah salah & separuh langkah salah & tidak ada atau acak\\\garis
B. Penjelasan & alasan tepat dengan istilah yang benar & tepat tetapi kurang lengkap & sebagian keliru & tidak ada atau disalin\\
\garisdasar
\end{tabular}}\par\vspace{6pt}
Skor soal = poin soal × skor level rata-rata aspek yang relevan.

\Bagian{Penutup}{Cara bersiap}
\begin{enumerate}
\item \textbf{Ulangi.} Modul Belajar 1 sampai 6 dan peta materi di slide P7.
\item \textbf{Latih.} Kerjakan empat contoh soal P8 dalam 120 menit tanpa bantuan.
\item \textbf{Siapkan.} Lembar catatan A4 tulisan tangan dan datang 15 menit lebih awal.
\end{enumerate}
\end{document}
